


This guide walks you through accessing UniZ, installing it on your device, and logging in for the first time. The whole process takes under five minutes.


\section{Instruction Step}

  
\section{Instruction Step}

    Navigate to the UniZ portal in any modern browser on your phone, tablet, or computer:

    \begin{lstlisting}[language=]
    https://uniz.rguktong.in
    
\end{lstlisting}

    You will see the UniZ home page with live campus announcements and a getting-started timeline.

  

  
\section{Instruction Step}

    UniZ is a Progressive Web App. Installing it gives you a faster, app-like experience without going through an app store.

    <Tabs>
      <Tab title="Android">
        1. Open `https://uniz.rguktong.in` in Chrome.
        2. Tap the **Install** button shown in the browser banner, or tap the three-dot menu and select **Add to Home Screen**.
        3. Confirm the installation prompt.

        UniZ now appears on your home screen and opens in full-screen mode.
      </Tab>
      <Tab title="iOS (iPhone / iPad)">
        1. Open `https://uniz.rguktong.in` in Safari.
        2. Tap the **Share** icon at the bottom of the screen.
        3. Scroll down and tap **Add to Home Screen**.
        4. Tap **Add** to confirm.

        <Note>
          iOS requires Safari for PWA installation. Chrome and other browsers on iOS do not support "Add to Home Screen" for PWAs.
        </Note>
      </Tab>
      <Tab title="Desktop (Windows / macOS / Linux)">
        1. Open `https://uniz.rguktong.in` in Chrome or Edge.
        2. Look for the install icon in the browser's address bar (a screen with a download arrow), or open the browser menu and select **Install UniZ**.
        3. Click **Install** to confirm.

        UniZ opens as a standalone desktop app.
      </Tab>
    </Tabs>

  

  
\section{Instruction Step}

    UniZ uses a single login endpoint for all roles. Where you log in depends on whether you are a student or a staff member.

    <Tabs>
      <Tab title="Students">
        1. From the portal home page, click **Get started** or navigate to the student sign-in page.
        2. Enter your university-issued student ID (for example, `O210008`) and your password.
        3. Click **Sign In**.

        **What you see after logging in:**
        - Your profile with name, branch, year, and campus status
        - Your grades and subject-wise attendance
        - Outpass and outing request forms
        - Your request history

        \begin{lstlisting}[language=json]
        {
          "username": "O210008",
          "password": "your_password"
        }
        
\end{lstlisting}

        <Tip>
          Your username is your RGUKT student roll number. If you have never logged in before, contact your administrator to get your initial credentials.
        </Tip>
      </Tab>
      <Tab title="Faculty & Staff">
        1. From the portal home page, type `admin` on your keyboard to reveal the admin sign-in page, or navigate directly to the admin portal.
        2. Enter your assigned username (for example, `dean_cse`, `warden_male`, `security_admin`) and your password.
        3. Click **Sign In**.

        **What you see after logging in:**
        - A dashboard tailored to your role
        - Pending outpass or outing requests (wardens, caretakers, SWO, dean, director)
        - Grade and attendance upload tools (faculty)
        - Gate check-in/check-out controls (security)
        - Banner and announcement management (webmaster)

        \begin{lstlisting}[language=json]
        {
          "username": "director",
          "password": "your_password"
        }
        
\end{lstlisting}

        <Note>
          All admin and faculty roles use the same login endpoint. The platform automatically loads the correct dashboard based on the role assigned to your account.
        </Note>
      </Tab>
    </Tabs>

  

  
\section{Instruction Step}

    Once you are logged in, here is what to do first for each role:

    | Role | First action |
    |------|-------------|
    | Student | Check your grades and attendance under the **Academics** section |
    | Caretaker / Warden | Open the **Requests** dashboard to see pending outpass approvals |
    | Faculty | Go to **Grades** and download the upload template for your batch |
    | Security | Open the **Gate** dashboard to see currently approved outpasses |
    | Dean / Director | Use **Search** to look up any student profile |
    | Webmaster | Go to **Banners** and create or publish a campus announcement |

  
</Steps>

\subsection{Forgot your password?}

If you cannot log in, you can reset your password using your registered email address.


\section{Instruction Step}

  
\section{Instruction Step}

    On the sign-in page, click **Forgot password** and enter your username. UniZ
    sends a one-time password to your registered email.
  
  
\section{Instruction Step}

    Enter the six-digit OTP from your email on the verification screen.
  
  
\section{Instruction Step}

    Enter and confirm your new password. You are then redirected to the sign-in
    page.
  
</Steps>

<Warning>
  OTPs expire after a short window. If you do not receive an email within a few
  minutes, check your spam folder or contact your system administrator.
</Warning>


\section{Deep Technical Analysis}
The underlying system architecture for this module is built upon the \textbf{UniZ Microservices Framework}. 
This design ensures that each functional unit (Auth, Profile, Academics) remains isolated, preventing cascading failures across the university ecosystem.

\subsection{Operational Hardening}
Deployments are managed via K3s (Kubernetes) and containerized using high-performance Alpine Linux Docker images. 
Resource management is critical; for instance, the documentation service itself is provisioned with 2GB of RAM to ensure the responsive rendering of complex documentation graphs and state-management components.

\subsection{Enterprise Security Topology}
Every request entering the UniZ gateway is subjected to an aggressive JWT validation layer. 
Permissions are granular, mapping directly onto the RBAC (Role-Based Access Control) matrix defined in the core system specifications.

\subsection{Data Governance}
Prisma-driven data access ensures type safety and predictable query performance. 
We utilize PostgreSQL connection pooling to handle concurrent requests from the student portal and administrative dashboards simultaneously.
